% !TEX root = ../notes.tex

\documentclass[../notes.tex]{subfiles}

\begin{document}

\section{April 22}
Today, we discuss Hecke operators.

\subsection{Hecke Operators}
Hecke operators come abstractly from our function space.
\begin{notation}
	Fix a reductive group $G$ over a global field $F$. For a level $K\subseteq G(\AA_F^\infty)$, we define $\mc A_G(K)$ to be the $K$-invariant automorphic forms. We also let $\mc H_{f,K}$ be the measures on $K\backslash G(\AA_F^\infty)/K$.
\end{notation}
\begin{remark}
	Note that $\mc A_G(K)$ embeds into $C^\infty(G(F)\backslash G(\AA_F)/K)$.
\end{remark}
\begin{remark}
	We may identify $\mc H_{f,K}$ with $C_c(K\backslash G(\AA_F^\infty)/K)$ by choosing $dg$ to have $\op{vol}(K;dg)=1$.
\end{remark}
For our Hecke algebras, note that whenever $G(\AA_F^\infty)$ acts on the right of a variety $X$, we receive a convolution left action of $\mc H_f$ on $C^\infty(X)$ by
\[(\mu*f)(x)\coloneqq\int_{G(\AA_F)}\mu(xg)h(g).\]
For example, if we extend from $\mc H_f$ to all measures so that we can plug in $h=\delta_g$, then $\delta_g*f$ is the right translation action.
\begin{remark}
	It is possible to write down some explicit formulae when we have fixed a level $K$: then the action of $h\,dg\in\mc H_{f,K}$ on $f\in\mc A_G(K)$ is given by
	\[(h*f)(x)=\sum_{g\in G(\AA_F^\infty)/K}f(xg)h(g).\]
	It is not hard to check that $h*f$ continues to be in $\mc A_G(K)$.
\end{remark}
\begin{example} \label{ex:hecke-action-correspondence}
	One convenience of Hecke operators is that they admit geometric description: let's describe the action of $1_{KgK}\in\mc H_{f,K}$. Namely, if $G(\AA_F^\infty)$ acts on a variety $X$, note that there is a correspondence as follows.
	% https://q.uiver.app/#q=WzAsNCxbMCwxLCJYL0siXSxbMCwwLCJYL1xcbGVmdChLXFxjYXAgZ0tnXnstMX1cXHJpZ2h0KSJdLFsxLDAsIlgvXFxsZWZ0KGdeey0xfUtnXFxjYXAgS1xccmlnaHQpIl0sWzEsMSwiWC9LIl0sWzEsMF0sWzEsMiwiUl9nIl0sWzIsM11d&macro_url=https%3A%2F%2Fraw.githubusercontent.com%2FdFoiler%2Fnotes%2Fmaster%2Fnir.tex
	\[\begin{tikzcd}[cramped]
		{X/\left(K\cap gKg^{-1}\right)} & {X/\left(g^{-1}Kg\cap K\right)} \\
		{X/K} & {X/K}
		\arrow["{R_g}", from=1-1, to=1-2]
		\arrow[from=1-1, to=2-1]
		\arrow[from=1-2, to=2-2]
	\end{tikzcd}\]
	It turns out that the Hecke operator is now computed by taking $f$ on $X/K$, pulling back along the projection, pulling back along $R_g$, and then pushing forward along the projection. This last projection is simply summing over the entire fiber, so we can see that we have given the correct operation.
\end{example}
While we're here, we interpret our Hecke operators in the function field setting, so let $X$ be a smooth projective integral curve over $\FF_q$ with function field $F$, and we use the group $\op{GL}_2$. Then our automorphic forms are simply some special functions $f$ on the groupoid of vector bundles of rank $2$.
\begin{notation}
	Fix a smooth projective integral curve $X$ over $\FF_q$ with function field $F$. For an effective divisor $D\subseteq X$, we define the Hecke operator $T_D$, defined as the indicator of the set
	\[\left\{g\in M_2(\widehat\OO_F):(\det g)=D\right\},\]
\end{notation}
\begin{remark}
	The element $\{T_D\}$ do not form a basis: we are missing the cosets $a\op{GL}_2(\widehat\OO)$ as $a\in\AA_F^\times$ varies.
\end{remark}
\begin{notation}
	For a divisor $D$ on $X$, we let $\delta_D$ be the indicator of $a\op{GL}_2(\widehat\OO)$, where $a\in\AA_F^\times$ represents the divisor $D$.
\end{notation}
\begin{remark}
	One can compute that $(\delta_D*f)(\mc E)=f(\mc E(D))$. Similarly, using the correspondence diagram, one finds
	\[(T_D*f)(\mc E)=\sum_{\substack{\mc E\subseteq\mc E'\\\op{div}(\mc E'/\mc E)=D}}f(\mc E').\]
\end{remark}
% Then for an effective divisor $D\subseteq X$, we have a 
% where the equality means that $\det g$ has matching support and valuations with the effective divisor $D$.

\subsection{Hecke Operators on Modular Forms}
For some intuition, we return to modular forms.
\begin{remark}
	In weight $k=2$, one can find the cusp forms $S_2(\Gamma)$ of level $\Gamma$ inside the cohomology $\mathrm H^1(\overline{\mc H/\Gamma};\CC)$ via the Hodge decomposition. Namely,
	\[\mathrm H^1(\overline{\mc H/\Gamma};\CC)=S_2(\Gamma)\oplus\overline{S_2(\Gamma)}.\]
\end{remark}
\begin{notation}
	For $\tau\in\mc H$, we let $E_\tau$ be the elliptic curve $\CC/(\ZZ+\tau\ZZ)$.
\end{notation}
\begin{lemma}
	The double quotient $\op{GL}_2(\ZZ)\backslash\op{GL}_2(\RR)/\CC^\times$ is in isomorphic with the double quotient of elliptic curves.
\end{lemma}
\begin{proof}
	We have two steps.
	\begin{enumerate}
		\item We claim that $\op{GL}_2(\ZZ)\backslash\op{GL}_2(\RR)/\CC^\times$ is in bijection with pairs $(\Lambda,J)$, where $\Lambda$ is a lattice of rank two, and $J$ is a complex structure on $\Lambda_\RR$. Indeed, $\op{GL}_2(\RR)$ simply classifies $2\times2$ matrices, which has a natural complex structure by identifying $\RR^2$ with $\CC$. Taking a quotient by $\CC^\times$ then produces lattices with a basis and some choice of complex structure $J$, and taking a further quotient by $\op{GL}_2(\ZZ)$  produces lattices with some choice of complex structure.
		\item We claim that pairs $(\Lambda,J)$ are in bijection with elliptic curve. Indeed, one can take $(\Lambda,J)$ to the complex manifold $\Lambda_\RR$ and then take quotient by the lattice $\Lambda$ to get an elliptic curve. Conversely, one can take an elliptic curve $E$ to the lattice $\mathrm H_1(E;\ZZ)$ equipped with the natural complex structure on $\mathrm H_1(E;\RR)$ canonically identified with $\op{Lie}E(\CC)$.
		\qedhere
	\end{enumerate}
\end{proof}
\begin{remark}
	Note that $\op{GL}_2(\RR)/\CC^\times$ has two components (corresponding to positive and negative determinant) which are then swapped by $\op{GL}_2(\ZZ)$. Thus, we could alternatively have worked with the double quotient $\op{SL}_2(\ZZ)\backslash\op{SL}_2(\RR)/\op{SO}_2(\RR)$ instead, which makes it clearer where the upper-half plane enters the picture.
\end{remark}
\begin{remark}
	It may be useful to notice that one can also interpret the quotient $\op{GL}_2(\ZZ)\backslash\op{GL}_2(\RR)$ in terms of elliptic curves. As above, if we undo the quotient by $\CC^\times$, this amounts to the data of a triple $(\Lambda,J,v)$ where $J$ is the complex structure on $\Lambda_\RR$, and $v\in\Lambda_\RR$ is some nonzero vector. The same argument now finds that such triples parameterize pairs $(E,v)$ of an elliptic curve $E$ together with a nonzero vector $v\in\op{Lie}E(\CC)$.
\end{remark}
The moral is that a modular form of level $k$ is now seen to be a function $f$ on pairs $(E,v)$ satisfying the functional equation
\[f(E,v)=c^kf(E,v)\]
for any $c\in\CC^\times$ (as well as some analytic conditions, such as being holomorphic and having moderate growth).

We are now ready to say something about Hecke operators. Let's attempt to write down a basis.
\begin{notation}
	For a positive integer $n$, we let $T_n$ be the indicator of the set
	\[\left\{g\in M_2(\widehat\ZZ):(\det g)=n\widehat\ZZ\right\}.\]
\end{notation}
\begin{example}
	For prime $p$, we find that
	\[\left\{g\in M_2(\widehat\ZZ):(\det g)=n\widehat\ZZ\right\}=\prod_{\ell\ne p}\op{GL}_2(\ZZ_\ell)\times\op{GL}_2(\ZZ_p)\begin{bmatrix}
		p \\ & 1
	\end{bmatrix}\op{GL}_2(\ZZ_p).\]
	Indeed, we are asking for trivial determinant away from $p$, and at $p$, we are asking for determinant $p$. The given description of the matrices with determinant $p$ in $\op{GL}_2(\QQ_p)$ follows from the Cartan decomposition.
\end{example}
\begin{example}
	For prime $p$, we similarly have
	\[T_{p^e}=\sum_{\substack{a+b=e \\a\le b}}1_{\op{GL}_2(\ZZ_p)\op{diag}\left(p^a,p^b\right)\op{GL}_2(\ZZ_p)}.\]
	Taking the product allows us to compute $T_n$ for any positive integer $n$.
\end{example}
\begin{remark}
	The elements $\{T_n\}$ are not a basis for the Hecke algebra: one needs to add in the commutative subalgebra of indicators of the double cosets $\op{GL}_2(\widehat\ZZ)\op{diag}(m,m)\op{GL}_2(\widehat\ZZ)$, where $m\in\QQ^+$. However, these extra operators are not so interesting: if $f$ is a modular form of weight $k$, then
	\[1_{m\op{GL}_2(\widehat\ZZ)}*f(E,v)=f(E,mv)=m^kf(E,v).\]
\end{remark}
Anyway, we now see that $\mc H_{f,K}$ is some commutative ring acting on our space $M_k(\op{GL}_2(\ZZ))$ of modular forms of weight $k$. It turns out that $M_k(\op{GL}_2(\ZZ))$ can be given a Hermitian structure, via the Peterson inner product
\[\langle f_1,f_2\rangle\coloneqq\int_{\mc H}f_1\overline{f_2}y^k\,\frac{dx\,dy}{y^2},\]
and one finds that the Hecke operators are self-adjoint.
\begin{remark}
	One can compute the adjoint of the Hecke action by $1_{KgK}$ on $M_k(\op{GL}_2(\ZZ))$ by using the correspondence diagram in \Cref{ex:hecke-action-correspondence}: the adjoint is found by going backward along the correspondence diagram, so it is basically $1_{Kg^{-1}\det(g)K}$, but one can show that $KgK=Kg^{-1}\det(g)K$ for $\op{GL}_2$.
\end{remark}
Anyway, the point is that the action by the Hecke algebra must diagonalize.
\begin{theorem}[Multiplicity one]
	The action of the Hecke algebra on the space of cusp forms for diagonalizes into one-dimensional eigenspaces.
\end{theorem}
\begin{definition}[eigenform]
	A cusp form $f$ is an \textit{eigenform} if and only if it is a simultaneous eigenvector for the entire Hecke algebra.
\end{definition}
\begin{remark}
	Suppose $f$ is an eigenform with $q$-expansion $f(q)=\sum_na_nq^n$. Then it turns out that $T_n*f$ is $a_n$ up to the scalar $a_1$. In particular, if $a_1=0$, then $f=0$; otherwise, we may normalize $f$ to have $a_1=1$, and then $T_n*f=a_n$.
\end{remark}
We are now ready to interpret our Hecke operators.
\begin{proposition}
	Fix a modular form $f$ of weight $k$ and full level. Then
	\[(T_n*f)(E,v)=\sum_{\substack{\pi\colon E'\to E\\\deg\pi=n}}f(E',v).\]
\end{proposition}
\begin{proof}
	Quickly, we note that the pair $(E',v)$ makes sense because an isogeny $\pi$ is automatically an isomorphism on the level of the Lie algebras.

	Anyway, the main point is to unwind everything to matrices. Given $(E,v)$, the data of an isogeny $E'\to E$ can be translated to lattices as follows: say $E=\CC/\Lambda$ and $E'=\CC/\Lambda'$, and then the isogeny of degree $n$ means that $\Lambda'\subseteq\Lambda$ with index $n$. However, once $\Lambda$ has been fixed, this collection is now in bijection with
	\[\{g\in M_2(\ZZ):\det g\in\{\pm n\}\}/\op{GL}_2(\ZZ),\]
	so one can find that everything unwinds as needed.
\end{proof}
Next time, we will say something about $L$-functions. Here is a preview for modular forms.
\begin{definition}
	Fix an eigenform $f$ with eigenvalues $\{a_n\}$. Then we define its $L$-function by
	\[L(s,f)\coloneqq\sum_{n=1}^\infty\frac{a_n}{n^s}.\]
	It is not obvious that this sum has any good properties.
\end{definition}
To say anything about this $L$-function, we relate it to the modular form.
\begin{lemma}
	Fix an eigenform $f$ with eigenvalues $\{a_n\}$. Then
	\[\int_{\RR^+}f(iy)y^s\,\frac{dy}y=(2\pi)^{-s}\Gamma(s)L(s,f).\]
\end{lemma}
\begin{proof}
	Direct calculation with $f(q)=\sum_na_nq^n$ (which holds after normalizing).
\end{proof}
\begin{theorem}
	Fix an eigenform $f$. Then $L(s,f)$ extends to an entire function and admits a functional equation.
\end{theorem}
\begin{proof}[Sketch]
	One can argue following Riemann: the functional equation $f(i/y)=f(-1/(iy))=(iy)^kf(iy)$ can be used to relate the integral for $L(s,f)$ with $L(1-s,f)$, so we receive a functional equation! Writing down the argument carefully enough gives the continuation as well.
\end{proof}
% As a preview, we note that an eigenform $f$ 
% However, it is not obvious that this integral has any good properties. Well, an explicit calculation with the $q$-expansion shows that
% \[\int_{\RR^+}f(iy)y^s\,\frac{dy}y=(2\pi)^{-s}\Gamma(s)L(s,f).\]
% Because $f$ is a cusp form, the requirement that $a_0=0$ implies that the integral automatically extends to $\op{Re}s>0$. For the rest of the continuation (and functional equation), we argue following Riemann: one can use the functional equation $f(i/y)=f(-1/(iy))=(iy)^kf(iy)$ to relate the integral for $L(s,f)$ with $L(1-s,f)$, so we receive a functional equation!

\end{document}